\chapter*{Roadmap}
\addcontentsline{toc}{chapter}{Roadmap}
\markboth{Roadmap}{Roadmap}

\section*{Why this course, and why this order}

The goal is a robot lawnmower that knows where it is on the lawn and builds
a map of it. A camera is the cheapest rich sensor we can put on it, so the
backbone of the course is \emph{visual} SLAM (Simultaneous Localisation And
Mapping). But a lawn is a notoriously hard scene for vision: the grass is a
repetitive texture, the light changes with clouds and time of day, shadows
move, and the few stable features (fences, trees, the house) are often far
away. Every serious outdoor mower therefore fuses vision with an inertial
measurement unit (IMU), wheel odometry and, when the sky is open, RTK-GNSS.
The course follows that path: first pure geometry and optimisation, then the
visual pipeline, then the back-end that makes it globally consistent, and
finally the fusion that makes it robust enough for the garden. The running
example is a concrete machine: a ride-on mower converted into a
robot. The \emph{On the lawn} boxes apply each chapter
to it; the rest of the course is not specific to it.

The course assumes linear algebra, multivariable calculus and probability
at the level of a first graduate course, some Python, and familiarity with
Bayesian filtering (the Kalman filter and, ideally, the particle filter).
Filtering is not re-taught: it is \emph{compared} with the
smoothing/optimisation view that dominates modern SLAM. For a refresher,
the companion notebooks at \url{https://github.com/CL-BZH/Bayesian_Filters}
derive the grid, Kalman, extended and unscented Kalman, and particle
filters, and compare them on a simulated robot mower: the same setting as
this course, seen through recursive filtering.

\section*{Structure}

\begin{description}[leftmargin=1.5em,style=nextline]
\item[Part I --- Foundations]
  \begin{enumerate}[label=\arabic*.,leftmargin=2em]
  \item \textbf{The SLAM problem.} Probabilistic formulation, the full joint
    posterior, MAP estimation as nonlinear least squares, factor graphs,
    filtering versus smoothing, observability and gauge freedom.
  \item \textbf{Rigid motion and Lie groups.} $\SO3$, $\SE3$, their Lie
    algebras, $\Exp$/$\Log$ in closed form, left and right Jacobians,
    adjoints, uncertainty on manifolds, perturbation calculus.
  \item \textbf{Cameras.} Pinhole projection, lens distortion, intrinsic
    calibration (Zhang's method), stereo and RGB-D models, the Jacobians of
    projection.
  \item \textbf{Nonlinear least squares on manifolds.} Gauss--Newton,
    Levenberg--Marquardt, trust regions, robust kernels (Huber, Cauchy) as
    iteratively reweighted least squares, optimisation with retractions.
  \end{enumerate}
\item[Part II --- The visual front-end]
  \begin{enumerate}[label=\arabic*.,leftmargin=2em,start=5]
  \item \textbf{Features and correspondence.} Corners (Harris, FAST),
    binary descriptors (ORB, BRIEF), matching, Lucas--Kanade optical flow,
    RANSAC and its probability analysis.
  \item \textbf{Two-view geometry.} Epipolar constraint, essential and
    fundamental matrices, eight- and five-point algorithms, decomposition
    of $E$, triangulation (linear and optimal), degenerate configurations
    (pure rotation, planar scenes --- the lawn!).
  \item \textbf{Visual odometry.} Perspective-$n$-Point, motion-only bundle
    adjustment, keyframes, a complete monocular/stereo VO pipeline, scale
    drift.
  \end{enumerate}
\item[Part III --- The back-end]
  \begin{enumerate}[label=\arabic*.,leftmargin=2em,start=8]
  \item \textbf{Bundle adjustment.} The full problem, sparsity structure,
    the Schur complement, gauge fixing, covariance recovery.
  \item \textbf{Factor graphs and pose graphs.} Pose-graph optimisation,
    sliding windows, marginalisation and its consistency issues
    (first-estimate Jacobians), incremental smoothing (iSAM2 ideas).
  \item \textbf{Loop closure and relocalisation.} Bag of visual words,
    place recognition, $\mathrm{Sim}(3)$ alignment for monocular scale,
    outlier-robust loop closure.
  \item \textbf{Maps for a mower.} Sparse versus dense maps, depth fusion,
    ground-plane extraction, a 2D occupancy/coverage map that plugs into
    the open-source \code{Lawnmower} coverage planners.
  \end{enumerate}
\item[Part IV --- Fusion for the garden]
  \begin{enumerate}[label=\arabic*.,leftmargin=2em,start=12]
  \item \textbf{The IMU.} Sensor model, biases and noise (Allan variance),
    kinematics on $\SO3$, on-manifold preintegration with full derivations.
  \item \textbf{Visual--inertial odometry.} Filtering (MSCKF) versus
    optimisation (sliding-window VIO), initialisation, observability.
  \item \textbf{Wheels and satellites.} Differential-drive odometry and
    slip, RTK-GNSS factors, frame alignment (ENU/local map), handling
    tree-canopy outages.
  \item \textbf{Capstone: a mower localisation stack.} Assembling the
    pieces, evaluation (ATE/RPE), a simulated garden, and the path to the
    real mower.
  \end{enumerate}
\end{description}

\section*{How each chapter is built}

Each chapter mixes four kinds of boxes. \textcolor{defblue}{Blue} boxes
are definitions, \textcolor{thmred}{red} boxes theorems,
\textcolor{propamber}{amber} boxes propositions and lemmas, all with full
proofs. \textcolor{wordsgreen}{Green} \emph{In words} boxes say what a
formula means and why we want it. Grey \emph{Practice} boxes point to code
in the companion repository, and \textcolor{lawngreen!70!black}{On the
lawn} boxes connect the material to the mower.

\section*{The companion code}

The Python package \code{vslam} grows chapter by chapter. Everything is
written from scratch with NumPy (and OpenCV only for image I/O and, later,
feature detection), so that every formula in these notes has a line of
code you can read. Every analytic Jacobian is checked against a numerical
one in the test suite. Where production systems do things differently, the
text points to the corresponding C++ components (Sophus, g2o, GTSAM,
Ceres, ORB-SLAM3, VINS-Fusion). No GPU is needed anywhere.

\begin{lstlisting}[language=bash,caption={Repository layout}]
vslam-course/
  course/          LaTeX sources of these notes
  code/vslam/      the Python package (lie.py, camera.py, ...)
  code/tests/      pytest suite, including Jacobian checks
  code/examples/   one script per practical
\end{lstlisting}

\section*{Data}

Synthetic data comes first: each practical generates its own scene, so
ground truth is exact and bugs are easy to localise. Real data enters in
Part~II with standard benchmarks: TUM RGB-D (indoor, RGB-D), EuRoC MAV
(stereo + IMU, ideal for Part~IV) and KITTI (outdoor stereo driving). For
the garden, the Rosario dataset (stereo, IMU, wheel odometry and RTK-GNSS
recorded by a wheeled robot in a soybean field) is the closest public
analogue to a lawn; the last step is to record sequences on the mower
itself.
